% Using InitechBase, Appendix B: Beginning and End of File (annex Sec 8.2
% UB, "Beginning/End of File and Filter Conditions").  Every row was
% observed through the shipped dot prompt (the examples below and those of
% Chapter 5); sources cmd/nav.c, cmd/query.c, cmd/mutate.c, cmd/report.c.
\chapter{Beginning and End of File}
\Gates{test-pubs-examples, test-interp-nav, test-samir-query}
\NoDictionaryHeads
The record pointer of a work area is always on a record, at the end of the
file after the last record, or, after a \cmd{SKIP} back from the first record,
at the beginning of the file.  \fn{EOF()} is true at the end of the file; there
\fn{RECNO()} is the number of records plus one.  \fn{BOF()} is true after a
movement back past the first record; \fn{RECNO()} is then 1.  The table gives
the position of the pointer after each command that moves it.

\noindent\hspace*{-\imHang}\begin{tabular}{@{}p{38mm}p{96mm}@{}}
\hline
\textsf{\textbf{Command}} & \textsf{\textbf{Record pointer afterwards}}\\
\hline
\cmd{USE} & The first record; with an index, the first in index order.\\
\cmd{GO TOP} & The first record, in index order and among the records the
  filter admits.  If no record is admitted, the end of the file.\\
\cmd{GO BOTTOM} & The last record, in the same sense.\\
\cmd{GO} \xp{n} & Record \xp{n}.\\
\cmd{SKIP} past the last record & The end of the file.  No message.\\
\cmd{SKIP} back past the first & The first record, with \fn{BOF()} true.  No
  message.\\
\cmd{LIST}, \cmd{DISPLAY ALL} & The end of the file.\\
\cmd{LIST NEXT} \xp{n}, \cmd{DISPLAY NEXT} \xp{n} & The last record listed.\\
\kw{WHILE} & The first record for which the condition is false.\\
\cmd{LOCATE}, \cmd{CONTINUE} & The record found, with \fn{FOUND()} true; if
  none, the end of the file.\\
\cmd{SEEK}, \cmd{FIND} & The first record with the key, with \fn{FOUND()} true;
  if none, the end of the file.\\
\cmd{REPORT FORM} & The end of the file.\\
\cmd{APPEND BLANK} & The new record.\\
\cmd{PACK} & The first record.\\
\cmd{ZAP} & The end of the file (record 1 of an empty database).\\
\hline
\end{tabular}

\section{Conditions that Hide Records}
Two settings make records invisible to the commands that move the record
pointer.  While SET DELETED is ON, records marked for deletion are passed over.
While SET FILTER is in force, records that do not meet the filter condition are
passed over.  \cmd{GO} \xp{n} still goes to record \xp{n} itself.  When every
record is hidden, \cmd{GO TOP} leaves the pointer at the end of the file:
\DotExample{examples/ub-filter}
